\section*{Abstract}

\ifshortabstract

\textbf{Background:} Randomised trials have established intravenous thrombolysis (IVT) and mechanical thrombectomy (MT) for ischaemic stroke, but they mostly report 90-day disability and have limited power to describe effects by stroke severity. We estimated effects on several outcomes by stroke severity and treatment strategy using national registry data.

\textbf{Methods:} We used target trial emulation in Sentinel Stroke National Audit Programme (SSNAP) data for patients with ischaemic stroke in England and Wales admitted in 2020--2025. Trials were defined by NIHSS band (1--5, 6--10, 11--15, 16--20, 21--42), treatment (IVT, MT or both), timing (within or after 4.5 hours of onset) and, for late treatment, perfusion imaging. Each treated patient was matched 1:1 to an untreated patient with identical pre-stroke disability, NIHSS band, age band and atrial fibrillation status. Outcomes at discharge were good functional outcome, modified Rankin Scale (mRS), in-hospital death, length of stay and discharge destination, compared using paired tests.

\textbf{Results:} Among 384,334 patients, 56,677 of 56,738 treated patients (99.9\%) were matched across 34 trial groups. Treatment was associated with more good outcomes in 28 groups, lower discharge mRS in 31, fewer deaths in 19 and a different mix of discharge destinations in 27, with more patients going home. Associations were small in mild stroke: with early IVT in NIHSS 1--5 there were 3.0 percentage points more good outcomes and slightly more deaths. The largest associations were with early IVT plus early MT in NIHSS 11--42: 23.3 to 34.2 percentage points more good outcomes and up to 21.3 percentage points fewer deaths. Associations were larger for early than late treatment. Length of stay was usually shorter, by up to 18.6 days.

\textbf{Conclusions:} Reperfusion treatment was associated with better discharge outcomes across most severity and treatment groups, most strongly for early combined treatment in moderate to severe stroke. Residual confounding is possible, so these estimates add detail to, but do not replace, trial evidence.

\else

% TODO: ~500-word abstract (used when \shortabstractfalse is set in main.tex).

\fi
